\documentclass[11pt,a4paper]{article}
\usepackage{../../common/preamble}

\begin{document}

\begin{center}
    {\LARGE \textbf{Базовая теория игр}}\\[0.2cm]
    \rule{\textwidth}{1pt}
\end{center}

\section{Математическая модель игры и понятийный аппарат}

В задачах 19--21 ЕГЭ исследуются \textbf{конечные антагонистические позиционные игры двух лиц с полной информацией и нулевой суммой}.

\begin{definitionbox}{Понятия теории позиционных игр}
\begin{itemize}
    \item \textbf{Позиция (состояние игры $s$):} набор параметров, полностью описывающий текущую игровую ситуацию (в игре с одной кучей --- число камней $s \in \mathbb{N}$, с двумя кучами --- пара $(a, b)$).
    \item \textbf{Правило перехода ($\text{moves}(s)$):} функция, возвращающая множество всех позиций, достижимых из $s$ ровно за один ход.
    \item \textbf{Дерево игры:} ориентированный граф, вершинами которого являются позиции, а дугами --- допустимые правилами ходы.
    \item \textbf{Выигрышная стратегия:} алгоритм выбора очередного хода, гарантирующий игроку победу при \textbf{любых} (в том числе абсолютно безошибочных) действиях оппонента.
    \item \textbf{Выигрышная позиция (Win):} позиция, в которой игрок, делающий очередной ход, обладает выигрышной стратегией.
    \item \textbf{Проигрышная позиция (Lose):} позиция, в которой игрок при любой своей игре неизбежно уступает, если соперник не совершает ошибок.
\end{itemize}
\end{definitionbox}

\section{Задача о 5 спичках}

Рассмотрим задачу для понимания логики обратной индукции:
\begin{quote}
\textit{На столе лежат 5 спичек. Два игрока по очереди убирают 1 или 2 спички. Побеждает тот, кто своим ходом оставит в кучке ровно 1 спичку.}
\end{quote}

Построим полное дерево игры:

\begin{center}
\begin{tikzpicture}[
    scale=0.95,
    level 1/.style={sibling distance=60mm},
    level 2/.style={sibling distance=30mm},
    level 3/.style={sibling distance=20mm},
    every node/.style={circle, draw, thick, minimum size=0.8cm, inner sep=1pt},
    edge from parent/.style={draw, thick, -Stealth}
]
    \node (root) [fill=black!5] {5}
        child {
            node (n4) [fill=black!10] {4}
            child {
                node (n3) [fill=black!10] {3}
                child {
                    node [fill=black!80, text=white] {1}
                    edge from parent node [left, draw=none] {\tiny $-2$}
                }
                edge from parent node [left, draw=none] {\tiny $-1$}
            }
            child {
                node (n2) [fill=black!10] {2}
                child {
                    node [fill=black!80, text=white] {1}
                    edge from parent node [right, draw=none] {\tiny $-1$}
                }
                edge from parent node [right, draw=none] {\tiny $-2$}
            }
            edge from parent node [above left, draw=none] {\footnotesize убрать 1}
        }
        child {
            node (n3_direct) [fill=black!10] {3}
            child {
                node [fill=black!80, text=white] {1}
                edge from parent node [right, draw=none] {\tiny $-2$}
            }
            edge from parent node [above right, draw=none] {\footnotesize убрать 2}
        };
\end{tikzpicture}
\end{center}

\subsection{Логический анализ дерева решений}
\begin{enumerate}
    \item Если Петя сходит $5 \to 3$ (уберет 2 спички), Ваня заберет 2 спички ($3 \to 1$) и \textbf{победит}.
    \item Если Петя сходит $5 \to 4$ (уберет 1 спичку), перед Ваней возникнет выбор:
    \begin{itemize}
        \item Ваня берет 1 спичку $\implies$ остается 3. Петя берет 2 спички и оставляет 1, побеждая.
        \item Ваня берет 2 спички $\implies$ остается 2. Петя берет 1 спичку и оставляет 1, побеждая.
    \end{itemize}
\end{enumerate}

\textbf{Вывод:} Позиция $4$ для игрока, делающего ход (для Вани), является \textbf{проигрышной}, так как \textbf{все} его возможные ходы ведут к выигрышу соперника. Следовательно, позиция $5$ для Пети является \textbf{выигрышной}, поскольку у него \textbf{существует ход} ($5 \to 4$), ставящий противника в заведомо проигрышную позицию.

\begin{definitionbox}{Кванторные критерии позиций}
\begin{itemize}
    \item Позиция $s$ является \textbf{выигрышной}, если существует хотя бы один ход в проигрышную позицию:
    $$s \in W \iff \exists h \in \text{moves}(s) : h \in L$$
    \item Позиция $s$ является \textbf{проигрышной}, если абсолютно все ходы ведут в выигрышные позиции:
    $$s \in L \iff \forall h \in \text{moves}(s) : h \in W$$
\end{itemize}
\end{definitionbox}

\section{Функция ранга позиции $R(s)$}

Метки $\{W, L\}$ не дают ответа на вопросы: \textit{\guillemotleft За сколько ходов побеждает игрок?\guillemotright} и \textit{\guillemotleft Гарантирована ли победа при любой игре соперника?\guillemotright}

Введем целочисленную функцию ранга позиции $R(s) \in \mathbb{Z}$, оценивающую состояние $s$ \textbf{относительно игрока, чья очередь делать ход}:
\begin{itemize}
    \item $R(s) = 0$ --- игра окончена (терминальное состояние, победа игрока, только что сделавшего ход)
    \item $R(s) > 0$ --- позиция выигрышная. Игрок побеждает ровно за $R(s)$ своих ходов.
    \item $R(s) < 0$ --- позиция проигрышная. Соперник гарантированно побеждает не более чем за $|R(s)|$ своих ходов.
\end{itemize}

Пусть из позиции $s$ доступны переходы в состояния $m \in \text{moves}(s)$, для каждого из которых вычислен ранг $R(m)$.  
Значение $R(m)$ оценивает позицию $m$ \textbf{с точки зрения нашего оппонента}, поскольку ход переходит к нему.
\begin{itemize}
    \item Если $R(m) > 0$, противник из позиции $m$ выигрывает. Нам такой ход невыгоден.
    \item Если $R(m) \le 0$, противник из позиции $m$ проигрывает (или игра уже завершилась при $R=0$). Для нас это \textbf{выигрышный переход}.
\end{itemize}

Следовательно, имеет смысл выделить из списка рангов $\texttt{r\_moves} \defeq \{R(m): m \in \text{moves}(s)\}$ подмножество неположительных рангов: $\texttt{neg} \defeq \{r \in \texttt{r\_moves if } : r \le 0\}$

Возможно 2 принципиально разных случая:

\subsubsection{$\texttt{neg} \ne \varnothing$}

Среди ходов есть хотя бы один переход в состояние, где противник терпит поражение. Текущий игрок выигрывает ($R(s) > 0$).

\textbf{Мотивация рационального игрока:} закончить партию как можно быстрее
\begin{itemize}
    \item Переход в $r = 0$ означает мгновенную победу (в 1 ход)
    \item Переход в $r = -1$ означает, что противник проиграет за 1 свой ход (мы побеждаем на 2-м ходу)
    \item Переход в $r = -2$ означает победу на 3-м ходу
\end{itemize}

Чтобы победить быстрее, нужно выбрать ход с наименьшим числом ответов противника: $\min |r| = \min (-r)$. Среди отрицательных чисел наименьший модуль имеет \textbf{максимальное число}:
$$\min \{1, 2\} \iff \max \{-1, -2\} = -1$$
Игрок выбирает ход $r^* = \max(\texttt{neg})$. Собственный ранг игрока:
$$R(s) = |r^*| + 1 = -r^* + 1 = -\max(\texttt{neg}) + 1$$

\subsubsection{$\texttt{neg} = \varnothing$}

Абсолютно все ходы ведут в $r > 0$. Куда бы ни сходил игрок, противник оказывается в выигрышной ситуации. Игрок проигрывает ($R(s) < 0$).

\textbf{Мотивация рационального игрока:} находясь в безнадежном положении, максимально затянуть сопротивление, рассчитывая на ошибку соперника
\begin{itemize}
    \item Ход в $r = 1$ позволяет сопернику победить сразу же (в 1 ход)
    \item Ход в $r = 2$ заставляет соперника делать сложную комбинацию из двух ходов
\end{itemize}

Игрок выбирает ход, требующий от противника наибольшего числа действий: $r^* = \max(\texttt{r\_moves})$. Так как игрок терпит поражение, его ранг отрицателен:
$$R(s) = -r^* = -\max(\texttt{r\_moves})$$

\section{Программный решатель}

\begin{taskbox}{Условие игры для одной кучи}
Два игрока, Петя и Ваня, играют в следующую игру. Перед игроками лежит куча камней. Игроки ходят по очереди, первый ход делает Петя.
За один ход игрок может добавить в кучу один камень или увеличить количество камней в куче в два раза. Для того чтобы делать ходы, у каждого игрока есть неограниченное количество камней.
Игра завершается в тот момент, когда количество камней в куче становится не менее 29. Победителем считается игрок, сделавший последний ход, т.е. первым получивший кучу, в которой будет 29 или больше камней.
В начальный момент в куче было $S$ камней, $1 \le S \le 28$. Будем говорить, что игрок имеет выигрышную стратегию, если он может выиграть при любых ходах противника.
\end{taskbox}

\subsection{Анализ задач}
\begin{enumerate}
    \item \textbf{№19.1:} Укажите такое значение $S$, при котором Петя не может выиграть за один ход, но при любом ходе Пети Ваня может выиграть своим первым ходом.

    Это соответствует определению проигрышной позиции ранга \textbf{$-1$}:
    $$R(s) == -1$$
    
    \item \textbf{№19.2:} Известно, что Ваня выиграл своим первым ходом после неудачного первого хода Пети. Укажите минимальное значение $S$, когда такая ситуация возможна.

    \textbf{Идея:} Ваня выигрывает своим 1-м ходом $\implies$ Петя сходил в состояние, где $R(m) = 1$. Значит, среди ходов Пети должен существовать ход в единицу: $\exists r \in \texttt{r\_moves} : r = 1$.
    
    При этом ход считается неудачным, если у Пети \textbf{был другой, более выгодный выбор}:
    \begin{itemize}
        \item либо Петя сам мог выиграть ($R(s) > 0$), но ошибся
        \item либо позиция проигрышная, но не все ходы вели к проигрышу в 1 ход ($\exists r \in \texttt{r\_moves} : r > 1$)
    \end{itemize}
    
    \item \textbf{№20:} Найдите два таких значения $S$, при которых у Пети есть выигрышная стратегия, причём одновременно выполняются два условия:
    \begin{itemize}
        \item Петя не может выиграть за один ход;
        \item Петя может выиграть своим вторым ходом независимо от того, как будет ходить Ваня.
    \end{itemize}

    Это позиция ранга \textbf{$+2$}:
    $$R(s) == 2$$

    \item \textbf{№21:} Найдите минимальное значение $S$, при котором одновременно выполняются два условия:
    \begin{itemize}
        \item у Вани есть выигрышная стратегия, позволяющая ему выиграть первым или вторым ходом при любой игре Пети;
        \item у Вани нет стратегии, которая позволит ему гарантированно выиграть первым ходом.
    \end{itemize}

    Это позиция ранга \textbf{$-2$}:
    $$R(s) == -2$$
\end{enumerate}

\subsection{Реализация базового решения}
\begin{verbatim}
T = 29

def moves(s: int) -> list[int]:
    """Возвращает позиции, доступные из данной."""
    return [s + 1, s * 2]

def R(s: int) -> int:
    """
    Как строится функция ранга позиции:
    1. R(s) = 0, если s >= T
    2.1. Получаем из s множество доступных позиций moves(s) 
         и строим множество их рангов R(moves(s))
    2.2. Выбираем из множества рангов все неположительные значения - neg. 
         Возможен один из двух вариантов:
         2.2.1. neg - пустое множество. Тогда R(s) = -max(R(moves(s)))
         2.2.2. neg - непустое множество. Тогда R(s) = -max(neg) + 1
    """
    if s >= T:
        return 0

    r_moves = [R(m) for m in moves(s)]
    neg = [r for r in r_moves if r <= 0]

    if neg:
        return -max(neg) + 1

    return -max(r_moves)

# 19.1
print(*(s for s in range(1, T) if R(s) == -1))

# 19.2
s_19_2 = []
for s in range(1, T):
    r_moves = [R(m) for m in moves(s)]
    if any(r == 1 for r in r_moves) and \
       (R(s) > 0 or not all(r == 1 for r in r_moves)):
        s_19_2.append(s)
print(*s_19_2)

# или некрасивый однострочник:
# print(*(s for s in range(1, T) if any(r == 1 for r in [R(m) for m in moves(s)]) 
#         and (R(s) > 0 or not all(r == 1 for r in [R(m) for m in moves(s)]))))

# 20
print(*(s for s in range(1, T) if R(s) == 2))

# 21
print(*(s for s in range(1, T) if R(s) == -2))
\end{verbatim}

\section{Практические приемы и усложнения}

\subsection{Оптимизация и кэширование}

Для больших диапазонов $S$ и игр с несколькими кучами прямое рекурсивное вычисление многократно пересчитывает одинаковые ветви, что приводит к сложности $\mathcal{O}(k^d)$. Для расчета за разумное время необходимо кэширование результатов.

\textbf{Вариант 1: Использование декоратора из стандартной библиотеки}
\begin{verbatim}
from functools import lru_cache

@lru_cache(None)
def R(s: int) -> int:
    ...
\end{verbatim}

\textbf{Вариант 2: Ручное кэширование через словарь}
\begin{verbatim}
results = {}

def R(s: int) -> int:
    if s in results:
        return results[s]
    
    if s >= T:
        return 0

    r_moves = [R(m) for m in moves(s)]
    neg = [r for r in r_moves if r <= 0]

    if neg:
        res = -max(neg) + 1
    else:
        res = -max(r_moves)
    
    results[s] = res
    return res
\end{verbatim}

\subsection{Ограничения на допустимые ходы}

Если правила запрещают определенные действия в зависимости от состояния, применяются два приема:

\textbf{Прием 1: Ветвление внутри функции \texttt{moves}}

\textit{Пример:} всегда разрешено $+2$, а также $\times 2$, если число нечетно, и $+1$, если число четно:
\begin{verbatim}
def moves(s: int) -> list[int]:
    m = [s + 2]
    if s % 2:
        m.append(s * 2)
    else:
        m.append(s + 1)
    return m
\end{verbatim}

\textbf{Прием 2: Обогащение состояния игры}

\textit{Пример:} игрокам запрещено повторять ход противника.

В состояние $s$ добавляется дополнительное поле: $0$ --- предыдущего хода не было, $1$ --- предыдущий ход был $+1$, $2$ --- предыдущий ход был $\times 2$:
\begin{verbatim}
def moves(s: tuple[int, int]) -> list[tuple[int, int]]:
    m = []
    if s[1] != 1:
        m.append((s[0] + 1, 1))
    if s[1] != 2:
        m.append((s[0] * 2, 2))
    return m

def R(s: tuple[int, int]) -> int:
    if s[0] >= T:
        return 0

    r_moves = [R(m) for m in moves(s)]
    neg = [r for r in r_moves if r <= 0]

    if neg:
        return -max(neg) + 1
    
    return -max(r_moves)

# Первичный вызов из начального состояния:
# R((s, 0))
\end{verbatim}

\subsection{Обобщение на игру с двумя кучами камней}

При добавлении второй кучи математическая структура алгоритма не меняется: состояние $s$ описывается парой $(a, b)$.

\begin{taskbox}{Игра с двумя кучами}
Два игрока, Петя и Ваня, играют в следующую игру. Перед игроками лежат две кучи камней. Игроки ходят по очереди, первый ход делает Петя.
За один ход игрок может:
\begin{itemize}
    \item добавить в одну из куч (по своему выбору) 4 камня;
    \item увеличить количество камней в одной из куч (по своему выбору) в 2 раза.
\end{itemize}
Например, пусть в одной куче 20 камней, а в другой 30 камней; такую позицию в игре обозначим (20, 30). Тогда за один ход можно получить любую из четырёх позиций: (24, 30), (20, 34), (40, 30), (20, 60).
Для того чтобы делать ходы, у каждого игрока есть неограниченное количество камней.
Игра завершается в тот момент, когда суммарное количество камней в двух кучах становится не менее 133. Победителем считается игрок, сделавший последний ход, то есть первым получивший такую игровую позицию, при которой в двух кучах суммарно 133 камня или больше.
В начальный момент в первой куче было 17 камней, во второй куче – $S$ камней; $1 \le S \le 115$.
Будем говорить, что игрок имеет выигрышную стратегию, если он может выиграть при любых ходах противника.
\end{taskbox}

\begin{verbatim}
from functools import lru_cache

T = 133
PILE_ONE = 17

def moves(state: tuple[int, int]) -> list[tuple[int, int]]:
    return [
        (state[0] + 4, state[1]), 
        (state[0] * 2, state[1]), 
        (state[0], state[1] + 4), 
        (state[0], state[1] * 2)
    ]

@lru_cache(None)
def R(state: tuple[int, int]) -> int:
    if state[0] + state[1] >= T:
        return 0

    r_moves = [R(m) for m in moves(state)]
    neg = [r for r in r_moves if r <= 0]

    if neg:
        return -max(neg) + 1

    return -max(r_moves)

# 19.1
print(*(s for s in range(1, T - PILE_ONE) if R((PILE_ONE, s)) == -1))

# 19.2
s_19_2 = []
for s in range(1, T - PILE_ONE):
    r_moves = [R(m) for m in moves((PILE_ONE, s))]
    if any(r == 1 for r in r_moves) and \
       (R((PILE_ONE, s)) > 0 or not all(r == 1 for r in r_moves)):
        s_19_2.append(s)
print(*s_19_2)

# 20
print(*(s for s in range(1, T - PILE_ONE) if R((PILE_ONE, s)) == 2))

# 21
print(*(s for s in range(1, T - PILE_ONE) if R((PILE_ONE, s)) == -2))
\end{verbatim}

\end{document}